\section{Conclusion}
\label{sec-conclude}

An ontology store is asked for patterns, hierarchies, and meaning.
\sys answers all three inside one disk-based, multi-threaded engine: a B+ tree
store of three permutations whose latch-free hits and optimistic writers scale to
eight threads; interval labels that make ``everything under $c$'' a few range
scans, which a plan rewrite uses to answer RDFS entailment with the answers of
materialization and none of its storage; vector indexes whose lists live in the
same tree; and a 0.85\,MB build for the browser.
It answers selective patterns and class instances faster than Oxigraph, Jena TDB2,
and SQLite and serves 3 to 7 times their lookups per second under eight threads.

\stitle{Limitations}.
\sys implements a \sparql subset without \texttt{OPTIONAL}, \texttt{UNION}, property
paths, or aggregates, its page-level recovery is exercised by a crash harness while
its durable database mode is a snapshot and logical log over the in-memory backend,
and it entails only \texttt{rdfs:subClassOf} over constant
classes; it is the slowest to load and uses much memory while loading; decoding an
id costs about $0.7\,\mu$s, which the cache hides on a warm store; labels and the
entailment snapshot are rebuilt, not maintained, after updates; and the vectors
come from one ontology and one small model, with synthetic instances.
